\documentclass[aspectratio=169]{beamer}

\usetheme{Madrid}
\usecolortheme{default}
\usefonttheme{professionalfonts}

\usepackage{amsmath,amssymb,bm}
\usepackage{physics} % for \ket,\bra,\dyad, optional
\usepackage{listings}
\usepackage{hyperref}

% ---------------------------
% Listings setup
% ---------------------------
\lstset{
  basicstyle=\ttfamily\small,
  columns=fullflexible,
  breaklines=true,
  frame=single,
  showstringspaces=false
}

\title{Bell States, Density Matrices, Partial Trace, and Measurements}
\author{}
\institute{}
\date{}

\begin{document}

\begin{frame}
  \titlepage
\end{frame}

% ==========================================================
\begin{frame}{1) The four Bell states (two qubits)}
Let $\{\ket{0},\ket{1}\}$ be the computational basis for a single qubit.
For two qubits, the computational basis is
$\{\ket{00},\ket{01},\ket{10},\ket{11}\}$.

\medskip
The \textbf{Bell states} are the four maximally entangled two-qubit states
\begin{align}
\ket{\Phi^{+}} &= \frac{1}{\sqrt{2}}\left(\ket{00}+\ket{11}\right), &
\ket{\Phi^{-}} &= \frac{1}{\sqrt{2}}\left(\ket{00}-\ket{11}\right),\\
\ket{\Psi^{+}} &= \frac{1}{\sqrt{2}}\left(\ket{01}+\ket{10}\right), &
\ket{\Psi^{-}} &= \frac{1}{\sqrt{2}}\left(\ket{01}-\ket{10}\right).
\end{align}

\medskip
They form an orthonormal basis of $\mathbb{C}^2\otimes\mathbb{C}^2$ and are
eigenstates of correlated observables (e.g.\ parity and certain two-qubit Pauli products).
\end{frame}

% ==========================================================
\begin{frame}{2) Density matrix for all four Bell states}
For a \textbf{pure} state $\ket{\psi}$, the density operator is
\[
\rho \;=\; \ket{\psi}\bra{\psi}.
\]

\medskip
Example:
\[
\rho_{\Phi^{+}}=\dyad{\Phi^{+}}
=\frac{1}{2}\big(\ket{00}\bra{00}+\ket{00}\bra{11}+\ket{11}\bra{00}+\ket{11}\bra{11}\big).
\]

\medskip
Similarly, expanding each Bell projector:
\begin{align}
\rho_{\Phi^{-}} &= \frac{1}{2}\big(\ket{00}\bra{00}-\ket{00}\bra{11}-\ket{11}\bra{00}+\ket{11}\bra{11}\big),\\
\rho_{\Psi^{+}} &= \frac{1}{2}\big(\ket{01}\bra{01}+\ket{01}\bra{10}+\ket{10}\bra{01}+\ket{10}\bra{10}\big),\\
\rho_{\Psi^{-}} &= \frac{1}{2}\big(\ket{01}\bra{01}-\ket{01}\bra{10}-\ket{10}\bra{01}+\ket{10}\bra{10}\big).
\end{align}
\end{frame}

% ==========================================================
\begin{frame}{2) Matrix forms in the computational basis}
Use ordering $\{\ket{00},\ket{01},\ket{10},\ket{11}\}$. Then

\[
\rho_{\Phi^{+}}=
\frac{1}{2}\begin{pmatrix}
1&0&0&1\\
0&0&0&0\\
0&0&0&0\\
1&0&0&1
\end{pmatrix},
\qquad
\rho_{\Phi^{-}}=
\frac{1}{2}\begin{pmatrix}
1&0&0&-1\\
0&0&0&0\\
0&0&0&0\\
-1&0&0&1
\end{pmatrix}
\]

\[
\rho_{\Psi^{+}}=
\frac{1}{2}\begin{pmatrix}
0&0&0&0\\
0&1&1&0\\
0&1&1&0\\
0&0&0&0
\end{pmatrix},
\qquad
\rho_{\Psi^{-}}=
\frac{1}{2}\begin{pmatrix}
0&0&0&0\\
0&1&-1&0\\
0&-1&1&0\\
0&0&0&0
\end{pmatrix}.
\]
\end{frame}

% ==========================================================
\begin{frame}{3) Proof that the Bell states form an orthonormal basis}
To show they form an orthonormal basis it suffices to prove:

\medskip
\textbf{(i) Normalization:} $\braket{\Phi^{\pm}}{\Phi^{\pm}}=1$ and $\braket{\Psi^{\pm}}{\Psi^{\pm}}=1$.
For example,
\[
\braket{\Phi^{+}}{\Phi^{+}}
=\frac{1}{2}\big(\braket{00}{00}+\braket{00}{11}+\braket{11}{00}+\braket{11}{11}\big)
=\frac{1}{2}(1+0+0+1)=1.
\]

\textbf{(ii) Orthogonality:} inner products between distinct Bell states vanish.
Example:
\[
\braket{\Phi^{+}}{\Phi^{-}}
=\frac{1}{2}\big(\braket{00}{00}-\braket{00}{11}+\braket{11}{00}-\braket{11}{11}\big)
=\frac{1}{2}(1-0+0-1)=0.
\]
Similarly, $\braket{\Phi^{\pm}}{\Psi^{\pm}}=0$ because $\ket{00},\ket{11}$ are orthogonal to $\ket{01},\ket{10}$.

\medskip
\textbf{(iii) Completeness:} There are 4 orthonormal vectors in a 4D space $\Rightarrow$ they form a basis.
\end{frame}

% ==========================================================
\begin{frame}{6) Pure and mixed states (definitions)}
\textbf{Density operator:} a quantum state can be represented by a positive semidefinite operator $\rho$
with $\Tr(\rho)=1$.

\medskip
\textbf{Pure state:}
\[
\rho \text{ is pure } \iff \rho=\ket{\psi}\bra{\psi} \iff \rho^2=\rho \iff \Tr(\rho^2)=1.
\]

\medskip
\textbf{Mixed state:} a probabilistic ensemble of pure states $\{\,(p_i,\ket{\psi_i})\,\}$,
\[
\rho=\sum_i p_i \ket{\psi_i}\bra{\psi_i},\qquad p_i\ge 0,\quad \sum_i p_i=1.
\]
Then typically $\rho^2\neq\rho$ and
\[
\Tr(\rho^2)<1.
\]

\medskip
\textbf{Example: maximally mixed single-qubit state} $\rho=\frac{I_2}{2}$ has $\Tr(\rho^2)=\frac{1}{2}$.
\end{frame}

% ==========================================================
\begin{frame}{4) Reduced density matrix via partial trace (full details)}
Let $\rho_{AB}$ be a two-qubit density matrix on $\mathcal{H}_A\otimes\mathcal{H}_B$.
The reduced state of qubit $A$ is
\[
\rho_A=\Tr_B(\rho_{AB})=\sum_{b\in\{0,1\}} \left(I_A\otimes \bra{b}\right)\rho_{AB}\left(I_A\otimes \ket{b}\right).
\]

\medskip
\textbf{Worked example for $\ket{\Phi^+}$:}
\[
\rho_{\Phi^+}=\frac{1}{2}\Big(\ket{00}\bra{00}+\ket{00}\bra{11}+\ket{11}\bra{00}+\ket{11}\bra{11}\Big).
\]

Take $\Tr_B$ term-by-term, using $\Tr_B\big(\ket{a b}\bra{c d}\big)=\braket{d}{b}\,\ket{a}\bra{c}$:

\begin{align}
\Tr_B(\ket{00}\bra{00}) &= \braket{0}{0}\ket{0}\bra{0}=\ket{0}\bra{0},\\
\Tr_B(\ket{00}\bra{11}) &= \braket{1}{0}\ket{0}\bra{1}=0,\\
\Tr_B(\ket{11}\bra{00}) &= \braket{0}{1}\ket{1}\bra{0}=0,\\
\Tr_B(\ket{11}\bra{11}) &= \braket{1}{1}\ket{1}\bra{1}=\ket{1}\bra{1}.
\end{align}

Therefore,
\[
\rho_A=\Tr_B(\rho_{\Phi^+})
=\frac{1}{2}\left(\ket{0}\bra{0}+\ket{1}\bra{1}\right)=\frac{I_2}{2}.
\]

\medskip
\textbf{Same result holds for all four Bell states:} each reduced state is maximally mixed.
\end{frame}

% ==========================================================
\begin{frame}{5) Measurement on one qubit (projective measurement)}
Consider measuring qubit $A$ in the computational basis $\{\ket{0},\ket{1}\}$.

\medskip
Projectors on the two-qubit space:
\[
P_0=(\ket{0}\bra{0})\otimes I,\qquad
P_1=(\ket{1}\bra{1})\otimes I.
\]

For a pre-measurement state $\rho_{AB}$,
\[
p(0)=\Tr(P_0\rho_{AB}),\qquad p(1)=\Tr(P_1\rho_{AB}), \qquad p(0)+p(1)=1.
\]

Post-measurement (collapsed) state conditioned on outcome $k\in\{0,1\}$:
\[
\rho_{AB}^{(k)}=\frac{P_k \rho_{AB} P_k}{p(k)}.
\]

\medskip
\textbf{Example:} for $\ket{\Phi^+}=(\ket{00}+\ket{11})/\sqrt{2}$, one finds
\[
p(0)=\frac{1}{2},\quad p(1)=\frac{1}{2},
\]
and
\[
\rho^{(0)}_{AB}=\ket{00}\bra{00},\qquad
\rho^{(1)}_{AB}=\ket{11}\bra{11}.
\]
Thus measuring qubit $A$ collapses qubit $B$ to the same outcome (perfect correlations).
\end{frame}

% ==========================================================
\begin{frame}{7) Pure-Python code: partial trace + von Neumann entropy}
\textbf{Goal:} compute $\rho_A=\Tr_B(\rho_{AB})$ and
\[
S(\rho_A)=-\Tr(\rho_A\log_2\rho_A)
= -\sum_{i=1}^{2}\lambda_i\log_2\lambda_i,
\]
where $\{\lambda_i\}$ are eigenvalues of the $2\times2$ matrix $\rho_A$.

\medskip
\textbf{No NumPy used:} we implement matrix ops with Python lists and compute eigenvalues analytically.
\end{frame}

\begin{frame}[fragile]{7) Code (standard Python libraries only)}
\begin{lstlisting}[language=Python]
import math
import cmath

def outer(v):
    """Outer product |v><v| for a statevector v (length 4)."""
    n = len(v)
    return [[v[i] * v[j].conjugate() for j in range(n)] for i in range(n)]

def partial_trace_B(rho):
    """
    Partial trace over qubit B for a 2-qubit density matrix rho (4x4).
    Basis ordering: |00>,|01>,|10>,|11>.
    Returns rho_A (2x2).
    """
    rhoA = [[0j, 0j],[0j, 0j]]
    # Indices: |a b> mapped to idx = 2*a + b
    for a in (0,1):
        for c in (0,1):
            s = 0j
            for b in (0,1):
                i = 2*a + b
                j = 2*c + b
                s += rho[i][j]
            rhoA[a][c] = s
    return rhoA

def eigvals_2x2(M):
    """Eigenvalues of a 2x2 matrix [[a,b],[c,d]] via quadratic formula."""
    a, b = M[0]
    c, d = M[1]
    tr = a + d
    det = a*d - b*c
    disc = tr*tr - 4*det
    root = cmath.sqrt(disc)
    lam1 = 0.5*(tr + root)
    lam2 = 0.5*(tr - root)
    # For density matrices these should be real (up to tiny imag parts)
    return [lam1.real, lam2.real]

def von_neumann_entropy(rhoA, eps=1e-15):
    """S(rhoA) = -Tr(rhoA log2 rhoA)."""
    lams = eigvals_2x2(rhoA)
    S = 0.0
    for lam in lams:
        lam = max(lam, 0.0)  # clip small negatives from numerical error
        if lam > eps:
            S -= lam * (math.log(lam, 2))
    return S

# --- Example: Bell state |Phi+> = (|00> + |11>)/sqrt(2)
inv_sqrt2 = 1.0 / math.sqrt(2.0)
phi_plus = [inv_sqrt2, 0j, 0j, inv_sqrt2]

rho = outer(phi_plus)
rhoA = partial_trace_B(rho)

print("rhoA =")
for row in rhoA:
    print(row)

print("Von Neumann entropy S(rhoA) =", von_neumann_entropy(rhoA))
\end{lstlisting}
\end{frame}

% ==========================================================
\begin{frame}{8) Pure-Python code: measurement on one qubit}
We implement projective measurement of qubit $A$ in $\{\ket{0},\ket{1}\}$.

\medskip
For a pure statevector $\ket{\psi}$, outcome probabilities are
\[
p(k)=\bra{\psi}P_k\ket{\psi},\quad P_k=(\ket{k}\bra{k})\otimes I,
\]
and the post-measurement state is
\[
\ket{\psi_k}=\frac{P_k\ket{\psi}}{\sqrt{p(k)}}.
\]

\medskip
Below: compute $p(0),p(1)$ and the conditional post-measurement statevectors.
\end{frame}

\begin{frame}[fragile]{8) Code (standard Python libraries only)}
\begin{lstlisting}[language=Python]
import math

def measure_qubit_A_statevector(psi, outcome):
    """
    Measure qubit A in computational basis on a 2-qubit statevector psi (len=4).
    outcome = 0 or 1.
    Returns (p, psi_post) where psi_post is normalized statevector.
    """
    # Projector keeps basis states with A=outcome:
    # indices 2*outcome + b for b in {0,1}
    keep = [2*outcome + 0, 2*outcome + 1]
    p = sum((psi[i].conjugate() * psi[i]).real for i in keep)
    if p <= 0.0:
        return 0.0, [0j,0j,0j,0j]
    norm = 1.0 / math.sqrt(p)
    psi_post = [0j,0j,0j,0j]
    for i in keep:
        psi_post[i] = psi[i] * norm
    return p, psi_post

# Example: |Phi+>
inv_sqrt2 = 1.0 / math.sqrt(2.0)
psi = [inv_sqrt2, 0j, 0j, inv_sqrt2]

p0, psi0 = measure_qubit_A_statevector(psi, 0)
p1, psi1 = measure_qubit_A_statevector(psi, 1)

print("p(0) =", p0, "post state =", psi0)
print("p(1) =", p1, "post state =", psi1)
\end{lstlisting}
\end{frame}

% ==========================================================
\begin{frame}{8) Qiskit code: Bell state + measurement on IBM hardware}
We create a Bell state, measure one qubit, and run on IBM Quantum using
Qiskit Runtime \texttt{SamplerV2}. The same circuit also works on simulators.

\medskip
Notes:
\begin{itemize}
\item The circuit prepares $\ket{\Phi^+}$: apply $H$ on qubit 0, then $\mathrm{CNOT}(0\to1)$.
\item Measure only qubit 0 (the ``one-qubit measurement'' request).
\item On real hardware, you need an IBM Quantum API token.
\end{itemize}
\end{frame}

\begin{frame}[fragile]{8) Qiskit Runtime (SamplerV2) example}
\begin{lstlisting}[language=Python]
from qiskit import QuantumCircuit
from qiskit.transpiler.preset_passmanagers import generate_preset_pass_manager

from qiskit_ibm_runtime import QiskitRuntimeService, SamplerV2 as Sampler

# 1) Authenticate (set your token once, e.g. via environment variable)
# service = QiskitRuntimeService(channel="ibm_quantum", token="YOUR_TOKEN")
service = QiskitRuntimeService()  # uses saved account configuration

# 2) Choose a backend (you can also specify a particular device name)
backend = service.least_busy(operational=True, simulator=False)

# 3) Build circuit: create Bell state |Phi+> and measure qubit 0 only
qc = QuantumCircuit(2, 1)
qc.h(0)
qc.cx(0, 1)
qc.measure(0, 0)

# 4) Transpile to the backend ISA
pm = generate_preset_pass_manager(backend=backend, optimization_level=1)
isa_circuit = pm.run(qc)

# 5) Run with SamplerV2
sampler = Sampler(mode=backend)
job = sampler.run([isa_circuit], shots=4096)
result = job.result()[0]

# Counts live under the classical register name; here it's "c"
counts = result.data.c.get_counts()
print("Backend:", backend.name)
print("Counts (measuring qubit 0):", counts)
\end{lstlisting}
\end{frame}

% ==========================================================
\begin{frame}{Quick checks / expected results}
For $\ket{\Phi^+}$ and measurement of qubit $A$ in the computational basis:

\begin{itemize}
\item $p(0)=p(1)=1/2$.
\item Conditional post-measurement state is $\ket{00}$ or $\ket{11}$.
\item Reduced state $\rho_A=I_2/2$ for any Bell state $\Rightarrow$
\[
S(\rho_A)=1 \text{ bit}.
\]
\end{itemize}

\medskip
These statements can be verified by the pure-Python scripts and by running the Qiskit circuit.
\end{frame}

\end{document}
